% paper/appendix_c_provenance_ledger.tex
% Full provenance-ledger mechanism description, moved from the former §3.7
% to Appendix C for the CFP's 10-page body limit. §3 now carries only a
% short summary ending "Full mechanism and the verification-bug episode
% described in Appendix C."

Every verdict is logged to an append-only SQLite ledger with a SHA-256 hash
chain: each row's \texttt{entry\_hash} is computed over the previous row's
hash plus a canonical serialization of that row's content-bearing fields, so
editing, deleting, or reordering any historical row breaks every subsequent
hash, detectable in $O(n)$ by \texttt{/ledger/verify}. The chain is
tamper-\emph{evident}, not tamper-\emph{proof}: an attacker with database
write access can recompute the whole chain, and defeating that requires an
external append-only anchor, left as future work. The hashing scheme is
schema-evolution-tolerant --- a field added after the ledger was already in
production only enters the hash for rows written after it existed, so old
rows continue to verify under the same code as new ones. We rely on this
directly: during this work we found and fixed a case where the
\emph{verification} query, not the ledger itself, omitted a newer column
from its recomputation, producing a false tamper report; correcting the
query restored a clean \texttt{intact: true} across all rows --- evidence
the mechanism is exercised in practice, not merely specified.
